\section{Factor systems}\label{sec:facsys}

Let $(\aA,G,\alpha)$ be a free C\Star-dynamical system.
Furthermore, for each $\sigma \in \Irrep(G)$ let $\hH_\sigma$ be a finite-dimensional Hilbert space and let $s(\sigma) \in \aA \tensor \End(V_\sigma,\hH_\sigma)$ an isometry satisfying $\alpha_g\bigl(s(\sigma)\bigr)=s(\sigma) \cdot \sigma_g$ for all $g \in G$ (cf.~\cite[Lemma~3.3]{SchWa17c}).
For $1 \in \Irrep(G)$ we choose $\hH_1 := \C$ and let $s(1) := \one_\aA$.
A~key feature of $(\aA,G,\alpha)$ is the factor system associated with the isometries $s(\sigma)$, $\sigma \in \Irrep(G)$, (see~\cite[Definition~4.1]{SchWa17c}), which we now recall for the convenience of the reader.
First, we put $\aB := \aA^G$.
Second, for expediency, we naturally extend $\sigma \mapsto \hH_\sigma$ and $\sigma \mapsto s(\sigma)$ to arbitrary finite-dimensional representations $\sigma$ of $G$ by taking the direct sum with respect to irreducible subrepresentations.
For each finite-dimensional representation $\sigma$ of $G$ we may then define the \Star-homomorphism
\begin{gather*}
	\gamma_\sigma\colon \ \aB \to \aB \otimes \End(\hH_\sigma),
	\qquad
	\gamma_\sigma(b) := s(\sigma) (b \tensor \one_{V_\sigma}) s(\sigma)^*,
\end{gather*}
and for each pair $(\sigma,\pi)$ of finite-dimensional representations of $G$ the element
\begin{gather*}
	\omega(\sigma,\pi) := s(\sigma) s(\pi) s(\sigma \tensor \pi)^* \in \aB \tensor \End(\hH_{\sigma \tensor \pi},\hH_\sigma \tensor \hH_\pi).
\end{gather*}
Then the following relations hold:
\begin{gather}
	\label{eq:ranges_sys}
	\omega(\sigma,\pi)^* \omega(\sigma,\pi) = \gamma_{\sigma \tensor \pi}(\one_\aB),
	\qquad
	\omega(\sigma,\pi) \omega(\sigma,\pi)^* = \gamma_\sigma (\gamma_\pi(\one_\aB) ),
	\\
	\label{eq:coaction_sys}
	\omega(\sigma, \pi) \gamma_{\sigma \tensor \pi}(b) = \gamma_\sigma( \gamma_\pi(b) ) \omega(\sigma, \pi),
	\\
	\label{eq:cocycle_sys}
	\omega(\sigma, \pi) \omega(\sigma \tensor \pi, \rho)
		= \gamma_\sigma ( \omega(\pi, \rho) ) \omega(\sigma, \pi \tensor \rho)
\end{gather}
for all finite-dimensional representations $\sigma$, $\pi$, $\rho$ of $G$ and $b \in \aB$ (see~\cite[Lemma~4.3]{SchWa20}).
The triple $(\hH, \gamma, \omega)$ of the above families is referred to as the \emph{factor system of $(\aA, G, \alpha)$} associated with $s(\sigma)$, $\sigma \in \Irrep(G)$, or simply as a factor system of $(\aA, G, \alpha)$ when no explicit reference to the isometries is needed.

We recall from~\cite{SchWa17c} that, for $\sigma \in \Irrep(G)$, the isotypic component~$A(\sigma)$ of $\aA$ can be written as
\begin{gather*}
	A(\sigma) = \bigl\{ \Tr ( y s(\sigma) )\colon y \in \aB \tensor \End(\hH_\sigma, V_\sigma) \bigr\}.
\end{gather*}
In fact, the map $y \mapsto \Tr ( y s(\sigma) )$ is bijective from $\aB \tensor \End(\hH_\sigma, V_\sigma) \gamma_\sigma(\one_\aB)$ to $A(\sigma)$.
The action~$\alpha$ on the isotypic component takes the form
\begin{gather}
	\alpha_g \bigl( \Tr ( ys(\sigma) ) \bigr) = \Tr\bigl( \sigma_g \cdot ys(\sigma) \bigr)\label{eq:A_0action}
\end{gather}
for all $g \in G$ and $y \in \aB \tensor \End(\hH_\sigma,V_\sigma)$.
The multiplication between isotypic components can be written as
	\begin{gather}
		\Tr \bigl( y_\sigma s(\sigma) \bigr) \cdot \Tr \bigl( y_\pi s(\pi) \bigr)
		=
		\Tr \bigl( y_\sigma \gamma_\sigma(y_\pi) \omega(\sigma, \pi) s(\sigma \tensor \pi) \bigr)\label{eq:A_0mul}
	\end{gather}
	for all $\sigma, \pi \in \Irrep(G)$, $y_\sigma \in \aB \tensor \End(\hH_\sigma, V_\sigma)$, and $y_\pi \in \aB \tensor \End(\hH_\pi, V_\pi)$.
The element on the right hand side is, in fact, a sum of elements in the isotypic components corresponding to subrepresentations of $\sigma \tensor \pi$.
Hence equation~\eqref{eq:A_0mul} uniquely determines the multiplication on the dense \Star-subalgebra~$\aA_f$ of~$\aA$.
The involution can be phrased in terms of the factor system, too (see~\cite[Section~5]{SchWa17c} or~\cite[Lemma~2.4]{Ne13}).
In order to see this, let us fix $\sigma \in \Irrep(G)$, choose an isometric intertwiner $v_0\colon \C \to V_{\bar\sigma} \tensor V_\sigma$, and write $w_0\colon \C \to \hH_{\bar\sigma \tensor \sigma}$ for the associated isometry.
Then the element $p_\sigma := (\dim \sigma)^2 w_0^{}v_0^*$ does not depend on the choice of the intertwiner.
It can be shown that there are element $\bar y_1, \dots, \bar y_n \in \aB \tensor \End(\hH_{\bar\sigma}, V_{\bar\sigma})$ and $y_1, \dots, y_n \in \aB \tensor \End(\hH_\sigma, V_\sigma)$ such that
\begin{gather}
	\label{eq:involution_decomposition}
	p_\sigma = \sum_{k=1}^n \bar y_k \gamma_{\bar\sigma}(y_k) \omega(\bar \sigma, \sigma)
\end{gather}
and that, for each $y \in \aB \tensor \End(\hH_\sigma, V_\sigma)$, the element
\begin{gather*}
	J(y) := \frac{1}{\dim \sigma}\sum_{k=1}^n \bar y_k \gamma_{\bar\sigma} \bigl( \Tr(y_k y^*) \bigr) \in \aB \tensor \End(\hH_{\bar\sigma}, V_{\bar \sigma})
\end{gather*}
does not depend on the choice of $\bar y_1, \dots, \bar y_n$ and $y_1, \dots, y_n$.
The involution on $\aA_f$ then reads as
\begin{gather*}
	\Tr \bigl( y s(\sigma) \bigr)^* = \Tr \bigl( J(y) s(\sigma)^* \bigr)
\end{gather*}
for all $y \in \aB \tensor \End(\hH_\sigma, V_\sigma)$.

Noteworthily, the notion of a factor system of $(\aA, G, \alpha)$ only depends on the fixed point algebra $\aB$ and the group $G$, which leads to the following definition:

\begin{Definition}\label{def:conjugancy}
	Let $\aB$ be a unital C\Star-algebra and let $G$ be a compact group.
	\begin{enumerate}\itemsep=0pt
	\item 		
		A \emph{factor system for $(\aB,G)$} is a triple $(\hH, \gamma, \omega)$ comprising families of Hilbert~spaces~$\hH_\sigma$, $\sigma \in \Irrep(G)$, \Star-homomorphisms $\gamma_\sigma\colon \aB \to \aB \tensor \End(\hH_\sigma)$, $\sigma \in \Irrep(G)$, and elements $\omega(\sigma, \pi)$ in $\aB \tensor \End(\hH_{\sigma \tensor \pi}, \hH_\sigma \tensor \hH_\pi)$, $\sigma, \pi\in \Irrep(G)$, satisfying equations~\eqref{eq:ranges_sys},~\eqref{eq:coaction_sys},~\eqref{eq:cocycle_sys}, as well as the normalization conditions $\hH_1 = \C$, $\gamma_1 = \id_\aB$, and $\omega(1, \sigma) = \gamma_\sigma(\one_\aB) = \omega(\sigma, 1)$ for all $\sigma \in \Irrep(G)$.
	\item
		Two factor systems $(\hH, \gamma, \omega)$ and $(\hH', \gamma', \omega')$ for $(\aB,G)$ are called \emph{conjugated} if there are partial isometries $v(\sigma) \in \aB \tensor \End(\hH_\sigma,\hH'_\sigma)$, $\sigma \in \Irrep(G)$, normalized to $v(1) = \one_\aB$, such that
		\begin{gather*}
			\Ad[v(\sigma)] \circ \gamma_\sigma = \gamma_\sigma',
			\qquad
			\Ad[v(\sigma)^*] \circ \gamma_\sigma' = \gamma_\sigma,		\\
			v(\sigma) \gamma_\sigma\bigl( v(\pi) \bigr) \omega(\sigma, \pi)
			= \omega'(\sigma, \pi) v(\sigma \tensor \pi)
		\end{gather*}
		for all $\sigma, \pi \in \Irrep(G)$.
		In this case we write $(\hH',\gamma', \omega') = v(\hH, \gamma, \omega)v^*$ or simply $(\hH,\gamma, \omega) \sim (\hH', \gamma', \omega')$ when no reference to the partial isometries is needed.
	\end{enumerate}
Note that, above, we have implicitly used functorial versions of the families $\hH_\sigma$, $\gamma_\sigma$, $v(\sigma)$, and $\omega(\sigma, \pi)$, $\sigma, \pi \in \Irrep(G)$.
\end{Definition}

By construction, each factor system of $(\aA,G,\alpha)$ is a factor system for $(\aB,G)$ and, by \cite[Lemma~4.3]{SchWa17c}, all factor systems of $(\aA, G, \alpha)$ are conjugated. In fact, given any unital C\Star\nobreakdash-algebra~$\aB$ and any compact group~$G$, we have shown in~\cite[Section~5]{SchWa17c} that the equivalence classes of free C\Star-dynamical systems $(\aA, G, \alpha)$ with fixed point algebra~$\aB$ are in 1-to-1 correspondence with conjugacy classes of factor systems of $(\aB,G)$.

\begin{Remark}\label{rem:K0_factor_sys}
For a factor system $(\hH, \gamma, \omega)$ of $(\aA,G,\alpha)$ equations~\eqref{eq:ranges_sys} and~\eqref{eq:coaction_sys} suggest to look at the $K$-theory of~$\aB$ and the induce positive group homomorphisms $K_0(\gamma_\sigma)\colon K_0(\aB) \to K_0(\aB)$ for a finite-dimensional representation~$\sigma$ of~$G$.
	Indeed, these maps only depend on the conjugacy class of the factor system and thus amount to invariants for $(\aA, G, \alpha)$.
	In fact, the mapping $\sigma \mapsto K_0(\gamma_\sigma)$ constitutes a nice functor:
	For direct sums $\sigma \oplus \pi$ of representation $\sigma$, $\pi$ of~$G$, we obviously have $K_0(\gamma_{\sigma \oplus \pi}) = K_0(\gamma_\sigma) + K_0(\gamma_\pi)$. Moreover, by equation~\eqref{eq:coaction_sys}, we have
	\begin{gather*}
		K_0(\gamma_{\sigma \tensor \pi}) = K_0(\gamma_\sigma) \circ K_0(\gamma_\pi)
	\end{gather*}
	for all finite-dimensional representations $\sigma$, $\pi$ of $G$.
\end{Remark}

\section{Lifting an automorphism}\label{sec:lift_automorphism}

Let $(\aA, G, \alpha)$ be a free C\Star-dynamical system with fixed point algebra $\aB$ and let $\beta$ be a \Star-auto\-mor\-phism of~$\aB$.
In this section we address the question whether $\beta$ can be lifted to a \Star-automorphism~$\hat \beta$ of~$\aA$ that commutes with all~ $\alpha_g$, $g \in G$.
In the affirmative case we say that \emph{$\beta$ lifts to $\aA$} and that $\hat \beta$ is a \emph{lift of~$\beta$}.

To phrase our result we note that $\Aut(\aB)$ acts on the factor systems for $(\aB,G)$.
For $\beta \in \Aut(\aB)$ and a factor system $(\hH, \gamma, \omega)$ for $(\aB,G)$ we may define a new factor system $\big(\hH, \gamma^\beta, \omega^\beta\big)$ for $(\aB,G)$ by putting
\begin{gather*}
	\gamma^\beta_\sigma
	:= \beta \circ \gamma_\sigma \circ \beta^{-1},
	\qquad
	\omega^\beta(\sigma, \pi)
	:= \beta \bigl( \omega(\sigma, \pi) \bigr)%\label{eq:beta_twist}
\end{gather*}
for all $\sigma, \pi \in \Irrep(G)$.
With this we give an answer to the above lifting problem by proving the following theorem:

\begin{Theorem}\label{thm:lift_conjugated_factor_sys}
	Let $(\aA, G, \alpha)$ be a free C\Star-dynamical system with fixed point algebra $\aB$ and let~$\beta$ be a \Star-automorphism of $\aB$.
	Then the following statements are equivalent:
	\begin{enumerate}\itemsep=0pt
	\item[$(a)$] $\beta$ lifts to $\aA$.
	\item[$(b)$] For any factor system $(\hH, \gamma, \omega)$ of $(\aA, G, \alpha)$ we have $(\hH, \gamma, \omega) \sim \big(\hH, \gamma^\beta, \omega^\beta\big)$.
	\end{enumerate}
\end{Theorem}
\begin{Remark}
	By the above theorem, a necessary condition for $\beta$ to lift to $\aA$ is that $K_0\big(\gamma_\sigma^\beta\big) = K_0(\gamma_\sigma)$ for all $\sigma \in \Irrep(G)$ or, equivalently, that $K_0(\gamma_\sigma)$ commutes with $K_0(\beta)$ for all $\sigma \in \Irrep(G)$ (cf.~Remark~\ref{rem:K0_factor_sys}).
	In particular, $K_0(\beta)$ is needs to fix the characteristic classes $[\gamma_\sigma(\one_\aB)] \in K_0(\aB)$, $\sigma \in \Irrep(G)$.
\end{Remark}

We split the proof of Theorem~\ref{thm:lift_conjugated_factor_sys} into a sequence of lemmas.
For a start we fix, for each $\sigma \in \Irrep(G)$, a finite-dimensional Hilbert space $\hH_\sigma$ and an isometry $s(\sigma)\in \aA \tensor \End(V_\sigma,\hH_\sigma)$ satisfying $\alpha_g (s(\sigma) )=s(\sigma) \cdot \sigma_g$ for all $g \in G$; for $1 \in \Irrep(G)$ we take $\hH_1 := \C$ and $s(1) := \one_\aA$.
As in Section~\ref{sec:facsys}, we write $(\hH, \gamma, \omega)$ for the associated factor system.
Our first result establishes the implication ``(a)$\Rightarrow$(b)'' of Theorem~\ref{thm:lift_conjugated_factor_sys}:

\begin{Lemma}\label{lem:conjugated}
	If $\beta$ lifts to $\aA$, then $(\hH, \gamma, \omega)$ and $\big(\hH, \gamma^\beta, \omega^\beta\big)$ are conjugated.
\end{Lemma}
\begin{proof}
	Let $\hat \beta$ be a lift of $\beta$.
	For each $\sigma \in \Irrep(G)$ we put $s^\beta(\sigma) := \hat \beta ( s(\sigma) )$ and note that $s^\beta(\sigma)$ is an isometry in $\aA \tensor \End(V_\sigma,\hH_\sigma)$ satisfying $\alpha_g\big(s^\beta(\sigma)\big)=s^\beta(\sigma) \cdot \sigma_g$ for all $g \in G$.
	Clearly, $s^\beta(1) = \one_\aA$.
	Furthermore, it is readily checked that the associated factor system is equal to $\big(\hH, \gamma^\beta, \omega^\beta\big)$.
	The~claim thus follows from the fact that all factor systems of $(\aA, G, \alpha)$ are conjugated.
\end{proof}

The task is now to prove the converse implication, ``(b)$\Rightarrow$(a)'', of Theorem~\ref{thm:lift_conjugated_factor_sys}.
For this purpose, we consider partial isometries $v(\sigma) \in \aB \tensor \End(\hH_\sigma)$, $\sigma \in \Irrep(G)$, normalized to $v(1) = \one_\aB$, such that $ \big(\hH, \gamma^\beta, \omega^\beta\big) = v (\hH, \gamma, \omega) v^*$.
Then for each $\sigma \in \Irrep(G)$ we obtain a well-defined map on the isotypic component $A(\sigma)$ by putting
\begin{gather}\label{eq:beta}
	\hat\beta_\sigma \bigl( \Tr (y s(\sigma)) \bigr) := \Tr \bigl( \beta(y) v(\sigma) s(\sigma) \bigr)
\end{gather}
for all $y\in \aB \tensor \End(\hH_\sigma, V_\sigma)$.
Taking direct sums gives a map~$\hat\beta$ on the dense \Star-subalgebra~$\aA_f$.
Due to the normalizations $v(1) = \one_\aB$ and $s(1) = \one_\aA$, we have $\hat \beta(b) = \beta(b)$ for all $b \in \aB$.
That~is, $\hat \beta$ extends $\beta$.
Furthermore, a few moments of thought show that $\hat \beta$ is bijective.
In fact, its inverse is given by the direct sum of the maps $\hat\beta^{-1}\colon A(\sigma) \to A(\sigma)$, $\sigma \in \Irrep(G)$, defined by
\[
\hat\beta^{-1}_\sigma \bigl(\Tr ( y s(\sigma)) \bigr)
	= \Tr \bigl( \beta^{-1} (y v(\sigma)^*) s(\sigma) \bigr)
\] for all $y \in \aB \tensor \End(\hH_\sigma, V_\sigma)$.
We proceed to establish further properties.

\begin{Lemma}\label{lem:prop_beta}
	The following assertions hold for the map $\hat \beta$ on $\aA_f$:
	\begin{enumerate}\itemsep=0pt
	\item[$1.$]
		$\hat\beta \circ \alpha_g = \alpha_g \circ \hat \beta$ for all $g \in G$.
	\item[$2.$]
		$\hat\beta$ is multiplicative.
	\item[$3.$]
		$\hat\beta$ is involutive.
	\end{enumerate}
\end{Lemma}
\begin{proof}1.~This is immediate from equations~\eqref{eq:A_0action} and~\eqref{eq:beta}.

2.~Let $\sigma, \pi \in \Irrep(G)$, let $x_\sigma = \Tr (y_\sigma s(\sigma)) \in A(\sigma)$ for some $y_\sigma \in \aB \tensor \End(\hH_\sigma, V_\sigma)$, and let $x_\pi = \Tr ( y_\pi s(\pi) ) \in A(\pi)$ for some $y_\pi \in \aB \tensor \End(\hH_\pi, V_\pi)$.
		Then
		\begin{align*}
			\hat\beta(x_\sigma) \cdot \hat\beta(x_\pi)
			&=
			\hat\beta \bigl( \Tr( y_\sigma s(\sigma) ) \bigr)
			\cdot \hat \beta \bigl( \Tr (y_\pi s(\pi) ) \bigr)
			\\
			&=
			\Tr \bigl( \beta(y_\sigma) v(\sigma) s(\sigma) \bigr)
			\cdot \id \tensor \Tr \bigl( \beta(y_\pi) v(\pi) s(\pi) \bigr)
			\\
			&=
			\Tr \bigl( \beta(y_\sigma) v(\sigma) \gamma_\sigma ( \beta(y_\pi) v(\pi) ) \omega(\sigma, \pi) s(\sigma \tensor \pi) \bigr).
		\end{align*}
		Using the conjugacy equations in Definition~\ref{def:conjugancy} now yields
		\begin{align*}
			\hat\beta(x_\sigma) \cdot \hat\beta(x_\pi)
			&=
			\Tr \bigl( \beta(y_\sigma) \gamma^\beta_\sigma ( \beta(y_\pi) ) v(\sigma) \gamma_\sigma ( v(\pi) ) \omega(\sigma, \pi) s(\sigma \tensor \pi) \bigr)
			\\
			&= \Tr \bigl( \beta ( y_\sigma \gamma_\sigma(y_\pi) ) \omega^\beta(\sigma, \pi) v(\sigma \tensor \pi) s(\sigma \tensor \pi) \bigr)
			\\
			&= \Tr \bigl( \beta ( y_\sigma \gamma_\sigma(y_\pi) \omega(\sigma, \pi) ) v(\sigma \tensor \pi) s(\sigma \tensor \pi) \bigr)
			=
			\hat\beta(x_\sigma \cdot x_\pi),
		\end{align*}
		which establishes that $\hat\beta$ is multiplicative.

3.~Let $\sigma \in \Irrep(G)$.
		To deal with the involution, we choose $\bar y_1, \dots \bar y_n \in \aB \tensor \End(\hH_{\bar\sigma}, V_{\bar\sigma})$ and $y_1, \dots, y_n \in \aB \tensor \End(\hH_\sigma, V_\sigma)$ satisfying equation~\eqref{eq:involution_decomposition} (cf.~Section~\ref{sec:facsys}).
		Now, let $x = \Tr ( y s(\sigma) )$ be in $A(\sigma)$ for some $y \in \aB \tensor \End(\hH_\sigma, V_\sigma)$.
		Then
		\begin{align*}
			\hat\beta(x^*)
			&= \hat\beta \bigl( \Tr ( J(y) s(\bar \sigma) ) \bigr)
			= \frac{1}{\dim \sigma} \sum_{k=1}^n \hat\beta \bigl( \Tr ( \bar y_k \gamma_{\bar\sigma} (\Tr(y_k y^*) ) s(\bar \sigma) ) \bigr)
			\\
			&=
			\frac{1}{\dim \sigma} \sum_{k=1}^n \Tr \bigl( \beta ( \bar y_k \gamma_{\bar \sigma} ( \Tr(y_k y^*) ) ) v(\bar \sigma) s(\bar\sigma) \bigr).
		\end{align*}
		It is easily checked that $\bar y_k^\beta := \beta(\bar y_k) v(\bar\sigma)$ and $y_k^\beta := \beta(y_k)v(\sigma)$, $1 \le k \le n$, also satisfy equation~\eqref{eq:involution_decomposition}.
		Since there is no loss of generality in assuming $y = y \gamma_\sigma(\one_\aB)$, we thus get
		\begin{align*}
			\hat\beta(x)^*
			&= \Tr ( \beta(y) v(\sigma) s(\sigma) )^*
			= \Tr \bigl( J ( \beta(y) v(\sigma) ) s(\bar\sigma) \bigr)
			\\
			&=
			\frac{1}{\dim \sigma} \sum_{k=1}^n \Tr \bigl( \beta(\bar y_k) v(\bar\sigma) \gamma_{\bar\sigma} ( \Tr ( \beta(y_k) v(\sigma) v(\sigma)^* \beta(y^*) ) ) s(\bar \sigma) \bigr)
			\\
			&=
			\frac{1}{\dim \sigma} \sum_{k=1}^n \Tr \bigl( \beta(\bar y_k) v(\bar\sigma) \gamma_{\bar\sigma} ( \beta \tensor \Tr( y_k y^*) ) s(\bar \sigma) \bigr).
		\end{align*}
		Invoking the conjugacy equations in Definition~\ref{def:conjugancy} finally gives
		\begin{align*}
			\hat\beta(x)^*
			&=
			\frac{1}{\dim \sigma} \sum_{k=1}^n \Tr \bigl( \beta(\bar y_k) \gamma^\beta_{\bar\sigma} ( \beta \tensor \Tr( y_k y^*) ) v(\bar\sigma) s(\bar \sigma) \bigr)
			\\
			&=
			\frac{1}{\dim \sigma} \sum_{k=1}^n \Tr \bigl( \beta ( \bar y_k) \gamma_{\bar\sigma}( \Tr( y_k y^*) ) v(\bar\sigma) s(\bar \sigma)\bigr)
			= \hat\beta(x^*),
		\end{align*}
		which shows that $\hat \beta$ is involutive.
\end{proof}

\begin{Lemma}\label{lem:beta_iso}
	We have $\rprod{\aB}{\hat\beta(x_1), \hat\beta(x_2)} = \beta\bigl( \rprod{\aB}{x_1, x_2} \bigr)$ for all $x_1, x_2 \in \aA_f$.
\end{Lemma}
\begin{proof}
	Since isotypic components are pairwise orthogonal, it suffices to show $\rprod{\aB}{\hat\beta(x_1), \hat\beta(x_2)} \allowbreak = \beta\bigl( \rprod{\aB}{x_1, x_2} \bigr)$ for all $x_1, x_2 \in A(\sigma)$ and $\sigma \in \Irrep(G)$.
	To this end, let $\sigma \in \Irrep(G)$ and let $x_1, x_2 \in A(\sigma)$.
	By Lemma~\ref{lem:prop_beta}, we obtain
	\begin{align*}
		\rprod{\aB}{\hat\beta(x_1), \hat\beta(x_2)}
		&=
		P_1\big( \hat\beta(x_1)^*\hat\beta(x_2) \big)
		=
		P_1\big( \hat\beta(x_1^*x_2) \big).
	\end{align*}
	We now decompose $x_1^*x_2$ as $\sum_{i \in I} \Tr( y_i s(\sigma_i) )$ for some mutually distinct representations $\sigma_i \in \Irrep(G)$ and $y_i \in \aB \tensor \End(\hH_{\sigma_i},V_{\sigma_i})$ and note that there is $i_0 \in I$ such that $\sigma_{i_0} = 1$, which is due to the fact that $\sigma \tensor \bar \sigma$ contains $1$ as a subrepresentation.
	Hence
	\begin{gather*}
		P_1\big( \hat\beta(x_1^*x_2) \big)
		=
		P_1\big( \beta(y_{i_0}) \big)
		=
		\beta\big( P_1(y_{i_0}) \big)
		=
		\beta\big( P_1(x_1^*x_2) \big)
		=
		\beta\bigl( \rprod{\aB}{x_1, x_2} \bigr).
		\tag*{\qed}
	\end{gather*}
	\renewcommand{\qed}{}
\end{proof}

By Lemma~\ref{lem:beta_iso}, the bijectivity of $\hat \beta$, and the fact that $\aA_f$ is dense in $L^2(\aA)$, we may extend~$\hat \beta$ to a unitary map, let's say, $U$ on $L^2(\aA)$.
Consequently, there is a \Star-automorphism on~$\aA$, for which we shall use the same letter $\hat \beta$ by a slight abuse of notation, such that
\begin{align*}
	\lambda\big( \hat \beta(x) \big) = U \lambda(x) U^*
\end{align*}
for all $x \in \aA$.
It is easily checked that $\hat \beta$ extends $\beta$ and commutes with all $\alpha_g$, $g\in G$.
Summarizing, we have shown the implication ``(b)$\Rightarrow$(a)'' of Theorem~\ref{thm:lift_conjugated_factor_sys}, which concludes the proof of this theorem.
